\documentclass{article}
\usepackage{mdwlist}
\usepackage[margin=2cm]{geometry}
\usepackage{graphicx}

\begin{document}

\begin{center}
\includegraphics[width=6cm]{ca/tex/syllabus/nsulogo} \\ \vspace{0.5cm}
College of Science and Health Professions \\
Department of Mathematics and Computer Science \\
Fall 2015
\end{center}

\begin{itemize}
\item[] \textbf{\Large MATH 1513-20139, College Algebra}

\item \textbf{Instructor:} Nathan Bloomfield, Ph.D.
\begin{itemize}
\item[] \textbf{Email:} \texttt{bloomfie@nsuok.edu}
\item[] \textbf{Office Location:} SC 252 
\item[] \textbf{Office Hours:} MWF 11--12, MW 8:30--9, 2--4
\item[] \textbf{Website:} \texttt{nbloomf.github.io}
\end{itemize}


\item \textbf{Course Delivery Mode:} Face-to-face


\item \textbf{Class Days and Times:} MWF 9--10 in SC 261


\item \textbf{Course Prerequisites and/or Corequisites:} Placement and enrollment in this course is based on ACT and/or CPT scores, or a combination of ACT and high school grades in mathematics.


\item \textbf{Catalog Description:} Linear and quadratic equations and inequalities; relations, functions, inverse functions; exponential and logarithmic functions; systems of equations; zeros of polynomials and determinants; permutations, combinations and the binomial theorem, as well as other selected topics.


\item \textbf{Course Purpose and Goals:} This course provides the mathematical development of the concept of a function, examines the basic types of functions (linear, quadratic, polynomial, exponential, logarithmic) and provides examples from a wide variety of applications ranging from business to the social sciences as well as in education and the natural sciences. The universal usefulness of this concept in most fields of study is why College Algebra is a required course for many majors.


\item \textbf{Course Topics:} We will cover material from the following sections in the textbook.
\begin{itemize*}
\item[] Chapter 1: Equations, Inequalities, and Modeling, Sections 1.1--1.7
\item[] Chapter 2: Functions and Graphs, Sections 2.1--2.6
\item[] Chapter 3: Polynomial and Rational Functions, Sections 3.1--3.5
\item[] Chapter 4: Exponential and Logarithmic Functions, Sections 4.1--4.4
\item[] Chapter 5: Systems of Equations and Inequalities, Sections 5.1--5.3
\item[] Chapter 7: The Conic Sections, Sections 7.1--7.3 (time permitting)
\end{itemize*}


\item \textbf{Student Learning Outcomes:} The student will be expected to achieve the following objectives.
\begin{enumerate*}
\item Communicate with appropriate mathematical symbols;
\item Solve equations and inequalities of the following types: linear, quadratic, absolute value, radical, polynomial, rational, exponential, logarithmic, containing 2 or 3 variables;
\item Define a function, identify and graph the various types of functions and perform algebraic operations on functions;
\item Model and solve real-world problems using linear, quadratic, polynomial, rational exponential and logarithmic functions;
\item Analyze how changing the parameters of functions affects the graphs of functions;
\item Convert between algebraic and graphical representations of rational, radical, absolute value, exponential and logarithmic relations; and
\item Compute efficiently using the appropriate method - calculator, paper/pencil algorithms, or mental calculation.
\end{enumerate*}


\item \textbf{General Education Learning Outcomes:} Specific educational objectives for the Quantitative Analysis category include the following.
\begin{enumerate*}
\item Solving problems using basic arithmetic and algebra;
\item Reasoning logically;
\item Communicating with symbols;
\item Drawing valid inferences from data presented in the form of a graph; and
\item Creatively applying known results to new situations.
\end{enumerate*}


\item \textbf{Instructional Methods:} This is a primarily lecture and demonstration-based course. However, many class meetings will include an activity assignment which will be worked on in small groups.


\item \textbf{Learning Outcome Assessment Methods:} Grades will be based on the following assignments.
\begin{itemize*}
\item[(60\%)] \textbf{Exams:} There will be 4 tests and a comprehensive final exam. The final will carry the same weight as two tests.
\item[(20\%)] \textbf{Activities:} We will have several in-class group activities.
\item[(10\%)] \textbf{Exam Reviews:} There will be a review assignment before each test and before the final exam.
\item[(10\%)] \textbf{``Daily'' Homework:} After most lectures there will be a very short homework assignment/quiz over that lecture's material.
\end{itemize*}
In each section I will also assign some homework problems from the textbook for practice. Although these will not be graded, it is strongly recommended that you work them.

The final grade will be the weighted average of the grades in each assignment category above. A final grade of 90 or better is an A; a grade in the interval $[80,90)$ is a B, et cetera. I reserve the right to adjust the cutoffs between letter grades downward at my discretion.

There will be \emph{one} extra credit opportunity: an algebra review assignment given out on the first day of the class. \textbf{There will be no other extra credit opportunities.} Also, students will be permitted to replace one chapter test grade with their final exam grade, provided the final exam grade is higher. \textbf{To take advantage of this, when the time comes, you must discuss it with me.}


\item \textbf{Instructional Materials.} 
\begin{itemize}
\item Required Text: \emph{College Algebra}, 5\textsuperscript{th} edition, by Mark Dugopolski (2011). The 4th edition of this text is also appropriate; see Dr. Bloomfield if you have questions.
\item A basic scientific calculator is recommended.
\end{itemize}

\item \textbf{Class and Instructor Policies:}
\begin{itemize}
\item \textbf{Attendance:} I do not give points for attendance. However, many class meetings will include graded assignments. It is to your benefit to come to class every day. If you are unable to come to class, plan to get notes and handouts from another student.

\item \textbf{Calculators:} Only a scientific, non-graphing, non-pretty-printing, dedicated calculator will be permitted on exams and quizzes. For instance, that means no TI-8x calculators or smartphones. \textbf{If you are not sure whether your calculator will be allowed on tests, ask me.} For all other non-exam and non-quiz work you are welcome to use whatever calculator you wish unless otherwise stated.

\item \textbf{Make-ups:} There will be no make-up tests without a good, documented reason. What counts as a ``good'' reason is up to me. If you know in advance that you will miss an exam (e.g. due to travel) let me know as soon as possible so we can schedule an alternative testing time.
\end{itemize}


\item \textbf{Academic Policies and Required Information:} Please go to 

\begin{center}
\texttt{http://offices.nsuok.edu/academicaffairs/SyllabiInformation.aspx}
\end{center}

for important information pertaining to:

\begin{itemize*}
\item Academic Misconduct
\item Americans with Disabilities Act (ADA) Compliance
\item Inclement Weather/Disaster Policy
\item Teach Act
\item Release of Confidential Information (FERPA)
\item Student Handbook
\item Textbook Information
\item Title IX
\end{itemize*}


\item \textbf{Class Calendar:} Test dates are to be determined. I will announce each test in class at least a week in advance.
\end{itemize}

\end{document}

